\documentclass{article}

\usepackage[utf8]{inputenc}
\usepackage[T1]{fontenc}
\usepackage[a4paper,margin=2.5cm]{geometry}
\usepackage{xcolor}
\usepackage{amsmath}
\usepackage{amssymb}
\usepackage{listings}

\lstset{
  basicstyle=\ttfamily\small,
  breaklines=true,
  frame=single,
  numbers=left,
  numberstyle=\tiny\color{gray},
  keywordstyle=\bfseries\color{blue!60!black},
  commentstyle=\itshape\color{green!40!black},
  stringstyle=\color{red!50!black},
  showstringspaces=false,
  tabsize=2
}

\title{Astrazione del Controllo nei Linguaggi di Programmazione}
\author{}
\date{}

\begin{document}

\maketitle

\section{Il Concetto di Astrazione}

\textbf{Astrarre} significa nascondere alcuni dettagli concreti per mettere in evidenza una struttura concettuale più generale. Nei linguaggi di programmazione l'astrazione è essenziale perché consente di descrivere problemi complessi in modo più semplice, modulare e manutenibile.

Si distinguono due grandi classi di meccanismi:
\begin{itemize}
  \item \textbf{astrazione del controllo}, che nasconde dati procedurali e schemi di esecuzione;
  \item \textbf{astrazione dei dati}, che consente di definire e usare tipi complessi indipendentemente dalla loro implementazione.
\end{itemize}

Questo documento riguarda l'\textbf{astrazione del controllo}: sottoprogrammi, passaggio dei parametri, funzioni di ordine superiore ed eccezioni.

\section{Sottoprogrammi}

Un \textbf{sottoprogramma}, \textbf{procedura} o \textbf{funzione} è una porzione di codice identificata da un nome, dotata di un ambiente locale e capace di comunicare con il resto del programma tramite:
\begin{itemize}
  \item parametri;
  \item valore di ritorno;
  \item ambiente non locale.
\end{itemize}

\subsection{Definizione e Chiamata}

\begin{lstlisting}
int foo(int n, int a) {
    int tmp = a;
    if (tmp == 0) return n;
    else return n + 1;
}

...
int x;
x = foo(3, 0);
x = foo(x + 1, 1);
\end{lstlisting}

La prima parte è la \textbf{definizione} della funzione; le ultime due istruzioni sono \textbf{chiamate}. La riga iniziale è detta \textbf{header}; il resto è il \textbf{corpo}.

\subsection{Parametri Formali e Attuali}

I \textbf{parametri formali} compaiono nella definizione della funzione e si comportano come nomi locali. I \textbf{parametri attuali} compaiono nella chiamata.

I parametri formali sono \textbf{variabili vincolate}: una rinomina coerente non modifica la semantica della funzione.

\begin{lstlisting}
int foo(int m, int b){
    int tmp = b;
    if (tmp == 0) return m;
    else return m + 1;
}
\end{lstlisting}

Questa funzione è semanticamente indistinguibile dalla precedente: cambiano i nomi formali, non la computazione.

\subsection{Valore di Ritorno}

Una funzione può restituire un valore al chiamante. Nei linguaggi della famiglia C, anche le procedure sono rappresentate linguisticamente come funzioni; quando il tipo di risultato è \texttt{void}, non viene restituito alcun valore significativo.

\subsection{Ambiente Non Locale}

Una funzione può comunicare con il resto del programma modificando variabili non locali. Questo meccanismo è meno strutturato rispetto a parametri e valori di ritorno e può indebolire l'astrazione funzionale.

\section{Astrazione Funzionale}

La \textbf{functional abstraction} separa ciò che il client deve conoscere da ciò che può essere nascosto:
\begin{itemize}
  \item il client conosce nome, parametri e tipo del risultato;
  \item il corpo della funzione è nascosto;
  \item l'implementazione può essere sostituita, se conserva la stessa semantica osservabile.
\end{itemize}

Si ha vera astrazione funzionale quando i client dipendono solo dall'header della funzione e non dal suo corpo.

\subsection{Variabili Statiche Locali}

In alcuni linguaggi una variabile locale può mantenere il proprio valore tra invocazioni diverse della stessa funzione. In C ciò si ottiene con il modificatore \texttt{static}.

\begin{lstlisting}[language=C]
int how_many_times(){
    static int count;
    /* C guarantees that static variables are
       initialised to zero by the first
       activation of the function */
    return count++;
}
\end{lstlisting}

Una variabile statica locale ha \textbf{lifetime} pari all'intera esecuzione del programma, ma \textbf{scope} limitato alla funzione. Offre quindi più astrazione rispetto a una variabile globale, perché non è modificabile dall'esterno.

\section{Passaggio dei Parametri}

La \textbf{disciplina di passaggio dei parametri} stabilisce come i parametri attuali sono associati ai parametri formali e quale semantica deriva da tale associazione.

Dal punto di vista della comunicazione, i parametri sono:
\begin{itemize}
  \item \textbf{input}, se comunicano dal chiamante al chiamato;
  \item \textbf{output}, se comunicano dal chiamato al chiamante;
  \item \textbf{input/output}, se la comunicazione è bidirezionale.
\end{itemize}

\subsection{Call by Value}

Il \textbf{call by value} implementa parametri di input. Alla chiamata, il parametro attuale viene valutato e il suo \textbf{r-value} viene copiato in una nuova variabile locale associata al parametro formale.

\begin{lstlisting}
int y = 1;

void foo(int x) {
    x = x + 1;
}

y = 1;
foo(y + 1);
// here y = 1
\end{lstlisting}

Durante l'esecuzione di \texttt{foo}, il parametro \texttt{x} assume valore iniziale $2$, viene incrementato a $3$, poi viene distrutto con il record di attivazione. La variabile \texttt{y} non cambia.

Proprietà principali:
\begin{itemize}
  \item l'attuale può essere una qualunque espressione;
  \item non c'è collegamento tra formale e attuale dopo la copia;
  \item non permette comunicazione dal chiamato al chiamante;
  \item è costoso per strutture dati grandi, perché richiede copia.
\end{itemize}

\subsection{Call by Reference}

Il \textbf{call by reference} implementa comunicazione input/output. Il parametro attuale deve avere un \textbf{l-value}; il formale diventa un alias dell'attuale.

\begin{lstlisting}
int y = 0;

void foo(reference int x) {
    x = x + 1;
}

y = 0;
foo(y);
// here y = 1
\end{lstlisting}

Durante la chiamata, \texttt{x} e \texttt{y} sono due nomi per la stessa locazione. Ogni modifica a \texttt{x} modifica \texttt{y}.

Il parametro attuale può anche essere un'espressione dotata di l-value:

\begin{lstlisting}
int[] V = new V[10];
int i = 0;

void foo(reference int x) {
    x = x + 1;
}

...
V[1] = 1;
foo(V[i + 1]);
// here V[1] = 2
\end{lstlisting}

Nell'implementazione con stack, nel record di attivazione del chiamato si memorizza l'indirizzo dell'attuale. Ogni accesso al formale richiede un'indirezione.

\subsection{Simulazione del Call by Reference in C}

C supporta solo il passaggio per valore, ma i puntatori permettono di simulare il passaggio per riferimento.

\begin{lstlisting}[language=C]
void swap(int *a, int *b) {
    int tmp = *a;
    *a = *b;
    *b = tmp;
}

int v1 = ...;
int v2 = ...;
swap(&v1, &v2);
\end{lstlisting}

I parametri \texttt{a} e \texttt{b} sono passati per valore, ma i loro valori sono indirizzi. Il dereferencing permette di modificare le variabili originali.

\subsection{Comunicazione Bidirezionale in Java}

Java passa i parametri per valore, ma per gli oggetti il valore copiato è una reference. È quindi possibile modificare lo stato dell'oggetto ricevuto, ma non sostituire la variabile del chiamante.

\begin{lstlisting}
class A{
    int v;
}

void swap(A a, A b) {
    int tmp = a.v;
    a.v = b.v;
    b.v = tmp;
}
\end{lstlisting}

Questo consente una forma di comunicazione bidirezionale sullo stato degli oggetti, ma non una vera simulazione generale del call by reference.

\subsection{Call by Constant}

Il \textbf{call by constant} è semanticamente equivalente al call by value, ma il formale non può essere modificato. È un parametro di input con vincolo statico di sola lettura.

Vantaggi:
\begin{itemize}
  \item documenta che il parametro è solo input;
  \item consente al compilatore di evitare copie costose usando una reference;
  \item preserva la semantica del call by value.
\end{itemize}

\subsection{Call by Result}

Il \textbf{call by result} è il duale del call by value. Implementa parametri solo output. Il formale è una variabile locale non inizializzata rispetto all'attuale; alla terminazione normale della procedura, il valore del formale viene copiato nella locazione dell'attuale.

\begin{lstlisting}
void foo(result int x) {
    x = 8;
}

...
int y = 1;
foo(y);
// here y is 8
\end{lstlisting}

Questioni semantiche rilevanti:
\begin{itemize}
  \item se ci sono più parametri result, occorre definire l'ordine delle copie finali;
  \item bisogna stabilire se l'l-value dell'attuale è calcolato alla chiamata o alla terminazione.
\end{itemize}

\subsection{Call by Value-Result}

Il \textbf{call by value-result} combina call by value e call by result. Alla chiamata si copia il valore dell'attuale nel formale; alla terminazione si copia il valore finale del formale nell'attuale.

\begin{lstlisting}
void foo(valueresult int x){
    x = x + 1;
}

...
y = 8;
foo(y);
// here y is 9
\end{lstlisting}

Il formale è una variabile locale. Per questo il call by value-result non coincide semanticamente con il call by reference in presenza di aliasing.

\begin{lstlisting}
void foo(reference/valueresult int x,
         reference/valueresult int y,
         reference int z){
    y = 2;
    x = 4;
    if (x == y) z = 1;
}

...
int a = 3;
int b = 0;
foo(a, a, b);
\end{lstlisting}

Se \texttt{x} e \texttt{y} sono value-result, non sono alias: il test può risultare falso e \texttt{b} non viene modificato. Se sono reference, \texttt{x} e \texttt{y} indicano la stessa locazione e il test è vero.

\subsection{Call by Name}

Il \textbf{call by name}, introdotto in ALGOL 60, è definito tramite la \textbf{copy rule}.

\textbf{Regola di copia.} Sia $f$ una funzione con un parametro formale $x$ e sia $a$ un'espressione compatibile con $x$. La chiamata $f(a)$ è semanticamente equivalente all'esecuzione del corpo di $f$ in cui ogni occorrenza di $x$ è sostituita da $a$, evitando cattura di variabili.

La sostituzione deve essere \textbf{capture-free}: l'attuale deve essere valutato nell'ambiente del chiamante, non nell'ambiente del chiamato.

\begin{lstlisting}
int x = 0;

int foo(name int y){
    int x = 2;
    return x + y;
}

...
int a = foo(x + 1);
\end{lstlisting}

Se l'espressione \texttt{x + 1} fosse valutata nell'ambiente del chiamato, la \texttt{x} locale catturerebbe la \texttt{x} del chiamante. Il call by name evita questa cattura associando all'espressione il suo ambiente originario.

\subsection{Effetti Collaterali nel Call by Name}

L'attuale viene rivalutato ogni volta che il formale è usato.

\begin{lstlisting}
int i = 2;

int fie(name int y) {
    return y + y;
}

...
int a = fie(i++);
// here i has value 4; a has value 5
\end{lstlisting}

La prima valutazione di \texttt{i++} produce $2$ e incrementa \texttt{i}; la seconda produce $3$ e incrementa di nuovo \texttt{i}. Il risultato è $5$.

\subsection{Call by Name e Value-Result non Coincidono}

\begin{lstlisting}
void fiefoo(valueresult/name int x,
            valueresult/name int y) {
    x = x + 1;
    y = 1;
}

...
int i = 1;
int[] A = new int[5];
A[1] = 4;
fiefoo(i, A[i]);
\end{lstlisting}

Con value-result, il riferimento a \texttt{A[i]} viene fissato rispetto al valore iniziale di \texttt{i}. Con call by name, \texttt{A[i]} viene rivalutato dopo la modifica di \texttt{i}; perciò può essere aggiornato un elemento diverso dell'array.

\subsection{Closure e Implementazione del Call by Name}

Per implementare il call by name, il chiamante passa una \textbf{closure}:
\[
\mathrm{closure} = (\mathrm{expression}, \mathrm{environment}).
\]

In una macchina astratta con stack, la closure è una coppia di puntatori:
\[
\mathrm{closure} = (\mathrm{code\ pointer}, \mathrm{static\ chain\ pointer}).
\]

Ogni accesso al parametro formale rivaluta l'espressione nell'ambiente memorizzato. Questo rende il call by name molto costoso.

\subsection{Jensen's Device}

Il call by name permette di passare una espressione e una variabile che compare nell'espressione, ottenendo procedure specializzabili al momento della chiamata.

\begin{lstlisting}
int summation(name int exp; name int i;
              int start; int stop) {
    int acc = 0;
    for (i = start, i <= stop, i++)
        acc = acc + exp;
    return acc;
}

int x = ...;
...
int y = summation(2*x*x - 2*x + 1, x, 1, 10);
\end{lstlisting}

La chiamata calcola:
\[
y = \sum_{x=1}^{10} (2x^2 - 2x + 1).
\]

\section{Funzioni di Ordine Superiore}

Una funzione è di \textbf{ordine superiore} se accetta funzioni come parametri o restituisce funzioni come risultato. Un linguaggio è detto \textbf{higher-order} se consente almeno una di queste due possibilità.

\subsection{Funzioni come Parametri}

\begin{lstlisting}
{
    int x = 1;

    int f(int y){
        return x + y;
    }

    void g(int h(int b)){
        int x = 2;
        return h(3) + x;
    }

    ...
    {
        int x = 4;
        int z = g(f);
    }
}
\end{lstlisting}

Il parametro \texttt{h} è una funzione. Quando \texttt{f} viene passata a \texttt{g}, occorre stabilire quale ambiente non locale usare per valutare \texttt{f}.

\subsection{Deep Binding e Shallow Binding}

Quando una funzione è passata come parametro, esistono due politiche di binding:
\begin{itemize}
  \item \textbf{deep binding}: l'ambiente è quello attivo quando si crea il legame tra parametro formale e funzione attuale;
  \item \textbf{shallow binding}: l'ambiente è quello attivo quando la funzione viene invocata tramite il parametro formale.
\end{itemize}

Nel frammento precedente:
\begin{itemize}
  \item con \textit{scope statico} e \textit{deep binding}, \texttt{h(3)} restituisce $4$ e \texttt{g} restituisce $6$;
  \item con \textit{scope dinamico} e \textit{deep binding}, \texttt{h(3)} restituisce $7$ e \texttt{g} restituisce $9$;
  \item con \textit{scope dinamico} e \textit{shallow binding}, \texttt{h(3)} restituisce $5$ e \texttt{g} restituisce $7$.
\end{itemize}

La \textbf{binding policy} è concettualmente indipendente dalla politica di scope, anche se i linguaggi con scope statico usano normalmente deep binding.

\subsection{Implementazione del Deep Binding}

Con deep binding, al parametro funzionale non basta associare il puntatore al codice: occorre associare anche l'ambiente non locale.

Il parametro funzionale è quindi rappresentato da una closure:
\[
\mathrm{closure}(f) = (\mathrm{code}(f), \mathrm{env}(f)).
\]

Quando il formale è invocato, la macchina astratta:
\begin{enumerate}
  \item trasferisce il controllo al codice memorizzato nella closure;
  \item inizializza il puntatore alla catena statica con l'ambiente memorizzato nella closure.
\end{enumerate}

\subsection{Binding Policy e Scope Statico}

Anche con scope statico, la binding policy è necessaria in presenza di ricorsione o mutua ricorsione, perché possono esistere più record di attivazione della stessa funzione.

\begin{lstlisting}
{
    void foo(int f(), int x){
        int fie(){
            return x;
        }
        int z;
        if (x == 0) z = f();
        else foo(fie, 0);
    }

    int g(){
        return 1;
    }

    foo(g, 1);
}
\end{lstlisting}

La funzione \texttt{fie} contiene una referenza non locale a \texttt{x}. Esistono però due attivazioni di \texttt{foo}. Con deep binding, l'ambiente usato è quello attivo quando \texttt{fie} viene associata al formale, quindi \texttt{z} riceve $1$. Con shallow binding riceverebbe $0$.

\subsection{Ingredienti dell'Ambiente}

Per determinare correttamente l'ambiente di valutazione in un linguaggio a blocchi servono:
\begin{itemize}
  \item regole di visibilità;
  \item eccezioni alle regole di visibilità;
  \item regole di scope;
  \item regole di passaggio dei parametri;
  \item binding policy.
\end{itemize}

\subsection{Funzioni come Risultati}

Se una funzione restituisce una funzione, il risultato non può essere rappresentato solo dal codice: deve includere anche l'ambiente in cui la funzione restituita dovrà essere valutata.

\begin{lstlisting}
{
    int x = 1;

    void->int F(){
        int g(){
            return x + 1;
        }
        return g;
    }

    void->int gg = F();
    int z = gg();
}
\end{lstlisting}

Il valore restituito da \texttt{F} è una closure:
\[
gg = (\mathrm{code}(g), \mathrm{env}(g)).
\]

\subsection{Upward Funarg Problem}

Se l'ambiente catturato dalla funzione restituita è locale alla funzione che termina, la disciplina stack dei record di attivazione non è più sufficiente.

\begin{lstlisting}
void->int F(){
    int x = 1;

    int g(){
        return x + 1;
    }

    return g;
}

void->int gg = F();
int z = gg();
\end{lstlisting}

Quando \texttt{F} termina, il suo record di attivazione verrebbe rimosso dallo stack, ma \texttt{gg} ha ancora bisogno dell'ambiente contenente \texttt{x}. La soluzione generale è allocare gli ambienti nel \textbf{heap} e deallocarli tramite \textbf{garbage collection}.

\section{Eccezioni}

Un'\textbf{eccezione} è un evento rilevato durante l'esecuzione che non deve o non può essere gestito nel normale flusso di controllo.

Può corrispondere a:
\begin{itemize}
  \item un errore semantico dinamico, come divisione per zero o overflow;
  \item una condizione anomala definita dal programmatore;
  \item una terminazione anticipata intenzionale di una computazione.
\end{itemize}

\subsection{Componenti di un Meccanismo di Eccezioni}

Un linguaggio con eccezioni deve specificare:
\begin{enumerate}
  \item quali eccezioni possono essere definite o gestite;
  \item come un'eccezione può essere sollevata;
  \item come un'eccezione può essere intercettata e gestita.
\end{enumerate}

In genere servono due costrutti:
\begin{itemize}
  \item un \textbf{protected block}, che delimita il codice in cui intercettare eccezioni;
  \item un \textbf{handler}, associato staticamente al protected block.
\end{itemize}

\subsection{Esempio di Gestione delle Eccezioni}

\begin{lstlisting}
class EmptyExcp extends Throwable {int x = 0;};

int average(int[] V) throws EmptyExcp(){
    if (length(V) == 0) throw new EmptyExcp();
    else {
        int s = 0;
        for (int i = 0, i < length(V), i++)
            s = s + V[i];
    }
    return s / length(V);
};

...
try {
    ...
    average(W);
    ...
}
catch (EmptyExcp e) {
    write("Array empty");
}
\end{lstlisting}

La clausola \texttt{throws} documenta che la funzione può terminare in modo anomalo. Il costrutto \texttt{throw} solleva l'eccezione. Il blocco \texttt{try} protegge una porzione di codice; \texttt{catch} contiene l'handler.

\subsection{Propagazione}

Quando un'eccezione viene sollevata:
\begin{enumerate}
  \item l'esecuzione del comando corrente viene interrotta;
  \item i blocchi attivi vengono abbandonati;
  \item la macchina astratta risale la catena dinamica;
  \item il controllo passa al primo handler compatibile;
  \item se non esiste un handler esplicito, interviene un handler di default.
\end{enumerate}

L'handler sostituisce la parte rimanente del protected block. Dopo l'esecuzione dell'handler, il controllo prosegue normalmente dopo il blocco di gestione.

\subsection{Handling con Terminazione e con Ripresa}

Il modello più comune è l'\textbf{handling with termination}: il costrutto in cui l'eccezione è stata rilevata termina. Alcuni linguaggi storici hanno usato \textbf{handling with resumption}, in cui l'handler può far riprendere l'esecuzione dal punto dell'errore; questo modello rende spesso gli errori più difficili da localizzare.

\subsection{Eccezioni e Catena Dinamica}

Le eccezioni propagano lungo la \textbf{catena dinamica}, anche se i nomi delle eccezioni seguono le regole statiche di scope.

\begin{lstlisting}
{
    void f() throws X{
        throw new X();
    }

    void g(int sw) throws X{
        if (sw == 0) { f(); }
        try {
            f();
        }
        catch (X e) {
            write("in g");
        }
    }

    ...
    try {
        g(1);
    }
    catch (X e) {
        write("in main");
    }
}
\end{lstlisting}

L'eccezione sollevata da \texttt{f} viene catturata dal più recente handler compatibile presente sullo stack di esecuzione.

\subsection{Uso Pragmatico delle Eccezioni}

Le eccezioni possono anche rendere più efficiente un algoritmo ordinario, interrompendo anticipatamente una computazione.

\begin{lstlisting}
type Node = {
    int key;
    Node FS;
    Node FD;
}

int mul(Node alb){
    if (alb == null){ return 1; }
    return alb.key * mul(alb.FS) * mul(alb.FD);
}
\end{lstlisting}

Se l'albero contiene uno zero, il prodotto è zero e non serve visitare il resto dell'albero.

\begin{lstlisting}
int mulAus(Node alb) throws Zero{
    if (alb == null){ return 1; }
    if (alb.key == 0) { throw new Zero(); }
    return alb.key * mulAus(alb.FS) * mulAus(alb.FD);
}

int mulEff(Node alb){
    try {
        return mulAus(alb);
    }
    catch (Zero e) {
        return 0;
    };
}
\end{lstlisting}

\texttt{mulEff} è un \textbf{wrapper}: esporta ai client un comportamento ordinario, mentre usa internamente un'eccezione per terminare in anticipo.

\section{Implementazione delle Eccezioni}

\subsection{Implementazione Basata sullo Stack}

Una soluzione semplice consiste nel memorizzare nel record di attivazione un puntatore all'handler associato al protected block, insieme al tipo di eccezione gestita.

Quando un'eccezione è sollevata:
\begin{enumerate}
  \item si cerca un handler nel record corrente;
  \item se non è presente, si ripristina lo stato della macchina;
  \item si rimuove il record dallo stack;
  \item si rilancia l'eccezione nel chiamante.
\end{enumerate}

Questa tecnica è concettualmente semplice ma introduce overhead anche nel caso normale, perché entrare e uscire da protected block richiede aggiornamenti dello stack.

\subsection{Implementazione con Tabella EH}

Una soluzione più efficiente sposta parte del lavoro a tempo di compilazione. Il compilatore costruisce una tabella nascosta $EH$ contenente, per ogni protected block, una coppia:
\[
(i_p, i_g)
\]
dove $i_p$ è l'indirizzo di inizio del protected block e $i_g$ è l'indirizzo di inizio dell'handler.

Quando un'eccezione è rilevata con program counter $pc$, la macchina astratta cerca una coppia tale che:
\[
i_p \leq pc \leq i_g
\]
con $i_p$ massimo tra quelli che soddisfano la condizione.

Se $EH$ è ordinata per $i_p$, la ricerca può essere binaria:
\[
T_{\mathrm{search}} = O(\log n)
\]
dove $n$ è il numero di handler registrati.

Questa tecnica costa di più quando l'eccezione è effettivamente sollevata, ma non costa nulla all'ingresso e all'uscita normale da un protected block.

\subsection{Eccezioni e Scope Statico}

I nomi delle eccezioni seguono lo \textbf{scope statico}. Questo può generare confusione quando due eccezioni hanno lo stesso nome in ambienti diversi.

\begin{lstlisting}
class X extends Exception {};

class P{
    void f() throws X{
        throw new X();
    }
}

class Q{
    class X extends Exception {};

    void g(){
        P p = new P();
        try {
            p.f();
        }
        catch (X e){
            System.out.println("in g");
        }
    }
}
\end{lstlisting}

Il metodo \texttt{f} solleva l'eccezione \texttt{X} dichiarata globalmente; il \texttt{catch} dentro \texttt{Q} intercetta invece \texttt{Q.X}. In Java ciò produce errori statici perché l'eccezione sollevata non è catturata e non è dichiarata nella clausola \texttt{throws} di \texttt{g}.

\section{Quadro Riassuntivo}

\begin{center}
\begin{tabular}{|p{3.2cm}|p{4.6cm}|p{5.6cm}|}
\hline
\textbf{Meccanismo} & \textbf{Comunicazione} & \textbf{Idea centrale} \\
\hline
Call by value & Input & Copia dell'r-value dell'attuale nel formale \\
\hline
Call by reference & Input/output & Il formale è alias dell'l-value dell'attuale \\
\hline
Call by constant & Input & Semantica per valore, implementazione anche per reference read-only \\
\hline
Call by result & Output & Copia finale dal formale all'attuale \\
\hline
Call by value-result & Input/output & Copia iniziale e copia finale, senza aliasing interno \\
\hline
Call by name & Input/output & Sostituzione capture-free dell'attuale e rivalutazione nell'ambiente del chiamante \\
\hline
Funzioni higher-order & Controllo come valore & Funzioni passate o restituite richiedono closure \\
\hline
Eccezioni & Salto controllato & Propagazione lungo la catena dinamica fino a un handler compatibile \\
\hline
\end{tabular}
\end{center}

\section{Sintesi Finale}

L'astrazione del controllo fornisce i meccanismi fondamentali per decomporre programmi complessi in componenti più piccoli, riusabili e semanticamente controllati. I \textbf{sottoprogrammi} realizzano astrazione funzionale; le \textbf{discipline di passaggio dei parametri} definiscono il canale di comunicazione tra chiamante e chiamato; le \textbf{funzioni di ordine superiore} trattano il controllo come valore; le \textbf{closure} catturano codice e ambiente; le \textbf{eccezioni} astraggono trasferimenti non locali di controllo.

Il tema comune è la separazione tra ciò che il programmatore vuole esprimere e i dettagli operativi della macchina astratta: record di attivazione, catene statiche e dinamiche, closure, handler e tabelle di eccezione sono i dispositivi implementativi che rendono possibile tale astrazione.

\end{document}
